\section{Discussion}

The cross-family result is asymmetric. The mapping \emph{questions} transfer well: incumbent region, exact trade witnesses, search/normal-form complexity, structural recognizability, and prospective prediction can all be posed for materially different compilers. The \emph{answers} need not share a form. TARE has a small exact support normal form but an evolving closed-form trade vocabulary. R6I collapses further to support one. SixLCU admits a strikingly low-order exact boundary. StabPrep does not admit an exact boundary in its frozen natural feature vocabulary at all.

This suggests that regime geometry should be treated like a phase-diagram methodology rather than a theorem schema. A useful map can contain a negative component: `no clean predicate in this vocabulary' is itself structural information, especially when mixed feature cells prove that more search over the same vocabulary cannot fix the classifier.

The support-bound results similarly distinguish validity from tightness. QG-6's production-inferred rank can guide theorem discovery, but QG-9 demonstrates that a semantics-derived finite bound may be loose by a large factor. Tightening a bound requires new transformations, not more confidence in the first certificate.

Objective dependence provides the final guardrail. A geometry belongs to a grammar--objective pair. QG-8 turns that guardrail into a theorem for one objective cone; outside the cone, a support-three control already exists. Comparisons across objectives should therefore report map transitions rather than assuming one canonical geometry.